\chapter{Conclusion}
\label{ch:conclusion}

\section{Summary of Contributions}
This thesis presented \textbf{MedPrepAI}, a comprehensive intelligent platform for virtual consultation, medical training, and clinical workflow, implemented across the entire clinical-lab software repository. The work extends far beyond a demonstration chatbot: it is a multi-module, deployable SaaS-shaped system suitable for institutional pilots.

Principal achievements include:
\begin{itemize}
  \item \textbf{Four-mode clinical AI curriculum}: Practice, Learning, Evaluation, and Shadow---each with distinct pedagogy, routes, entitlements, and persistence (Figure~\ref{fig:mode-pedagogy}).
  \item \textbf{Full LMS}: Hierarchical content, timed assessments, mock exams, rich explanations, DOCX/Markdown ingestion.
  \item \textbf{Institutional operations}: Faculty analytics, assignments, student comparison, domain-linked cohorts.
  \item \textbf{Governed commerce}: Stripe billing, subscription packages, modern entitlement types (BOOLEAN, SET, NUMBER\_LIMIT).
  \item \textbf{Collaboration layer}: Chat, notifications, discussions, study groups, achievements.
  \item \textbf{Engineering excellence}: 25+ NestJS modules, 20 Prisma schema files, documented deployment, demo seed covering all subsystems.
\end{itemize}

The requirements traceability matrix (Table~\ref{tab:traceability}) shows all FR-01--FR-10 passed scenario-based evaluation. Non-functional requirements NFR-01--NFR-06 were satisfied through guards, indexing, modular schemas, audit capability, deployment documentation, and educational LLM disclaimers.

\section{Comparison with Initial Goals}
At project inception, goals focused on porting UWorld-style UI and adding login. The final system substantially exceeded that scope by delivering MedPrep AI modes, faculty workspace, payments, entitlements, and community modules---reflecting iterative expansion aligned with supervisor feedback and institutional user stories.

\section{Limitations}
\begin{enumerate}
  \item \textbf{Clinical validity}: AI patients are simulations; outcomes were not validated against practicing clinicians' gold-standard ratings.
  \item \textbf{LLM variability}: Gemini model updates may alter dialogue style; continuous prompt regression is required.
  \item \textbf{Evaluation scale}: No randomized controlled trial with medical students; results are scenario- and demo-seed-based.
  \item \textbf{Regulatory scope}: Not FDA-cleared; not for real patient care.
  \item \textbf{Automated test coverage}: E2E test suite remains future work; manual QA dominated capstone timeline.
  \item \textbf{Legacy naming}: Early repository templates referenced unrelated domains (rides); documentation now centers LMS/MedPrep paths.
\end{enumerate}

\section{Lessons Learned}
\begin{itemize}
  \item \textbf{Separate LLM roles early}: Patient, tutor, and evaluator prompts must not share unconstrained system strings.
  \item \textbf{Persist early}: Moving MedPrep JSON files into Prisma (\texttt{migrate-medprep-json}) simplified faculty analytics.
  \item \textbf{Entitlements as data}: SET-type MedPrep mode lists avoid redeploying when marketing changes packages.
  \item \textbf{Cap metadata size}: Shadow interventions required explicit caps to keep MySQL JSON rows performant.
\end{itemize}

\section{Future Work}
\begin{itemize}
  \item Multilingual patients and rubrics; Urdu/English code-switching for regional institutions.
  \item Retrieval-augmented generation over licensed medical corpora with citation cards in Learning mode.
  \item On-premise LLM deployment for hospitals prohibiting cloud PHI (even if synthetic cases).
  \item Automated multimodal proctoring in Evaluation mode.
  \item Competency graphs mapping encounters to ACGME milestones.
  \item FHIR sandbox integration for documentation export training.
  \item Playwright E2E suite in CI; load testing MedPrep APIs at 500 concurrent sessions.
  \item Peer-reviewed usability study with Year-3--5 medical students.
\end{itemize}

\section{Closing Remarks}
MedPrepAI shows that the next generation of medical education software must unify AI consultation, assessment science, and institutional workflow. The clinical-lab repository provides an extensible foundation---students Zayan Rashid Rana and Muhammad Ahad Siddique implemented not a feature demo but a \emph{platform}. We thank Dr.\ Kamran Safdar and Ahmad Hassan Chaudhry for guidance, and PIEAS for the opportunity to deliver this capstone contribution to healthcare informatics education.
